\subsection{Academy}

Academy er et kortspil for $2$ - $6$ personer, der især blev populært under corona hvor man kunne spille det online på \href{https://academy.beer/}{academy.beer}. Spillet handler om at man trækker kort og så drikker det der står på kortet.

Historien siger at Academy blev lavet da der var nogle på \TKET der blev trætte af at folk snød i mængde de skulle drikke i diverse drukspil. Derfor lavede man Academy! 

Spillet fungerer ved at man tager et dæk kort og en kulør per spiller\footnote{Man fjerne de kulører der ikke skal bruges.} og så får den første det øverste kort fra bunken og skal drikker det antal slurke der står på kortet. Så får den næste spiller det næste kort fra bunken og drikker det antal slurke der står på det kort osv... Når det så bliver den førstes tur igen, må de først få et kort når de har drukket det antal slurke som der stod på det forrige kort. Så siger man at man er på ``niveau", hvor et nyt kort bliver trukket. Sådan fortsætter det indtil dækket er tomt.


Der er et par regler i spillet:
\begin{itemize}
	\item Inden spillets start bliver man enige om, hvor mange tårer der er i én 33cl øl. Standard er 14 tårer.
	\item Kortene har værdi: Es: 14, Konge: 13, Dronning: 12, ... 3'er: 3, 2'er: 2. Dvs. hvis man trækker et Es så skal man bunde en ny øl.
	\item Man skal have drukket det antal slurke der står på kortet inden man må trække det næste kort. Man får det næste kort ved at sige man er på ``niveau".
	\item Man må ikke trække sit eget kort, det skal gøres af en medspiller, for at undgå snyd.
	\item Man må gerne drikke foran, men maximalt én øl. Alt efter det er interessedruk
	\item Bund reder ikke.
\end{itemize}


Det online Academy fungere meget som det analoge, der er blandt andet blevet tilføjet to kulører, nemlig Heineken og Carlsberg logoet. Ellers er der kun ændringen at man nu trykker på mellemrumstasten, også kaldet ``Den Lange Tast", får at få et nyt kort. Det er stadig en medspiller der skal trykke på den lange tast for at undgå snyd.

 I online Academy har hver person også en profil de spiller på, hvilket gør det muligt at lave statistik på de forskellige spil og spillere. Inde på \href{https://academy.beer/}{academy.beer} kan man finde statistik omkring gennemsnitlig bundetider, gennemsnitlig spillænge, total antal esser en person har fået og meget mere. Inde på hjemmesiden står der også mere omkring historien om online Academy.
 
 Det online Academy har også mulighed for at tilknytte en Discord konto til sin profil. Hvis man har gjort dette får man besked via Discord, hvornår det er ens tur og hvor meget man skal drikke. Dette kan være meget smart, hvis man spiller med folk der render rundt hele tiden.